\subsection*{Order Relations}
\label{subsec:order-relations}

\begin{definitionbox}{Definition (Partial Order)}
\begin{definition}[Partial Order]
\label{def:algebraic-partial-order}
A binary relation $R$ on $A$ is a \textbf{partial order} if it is reflexive,
antisymmetric, and transitive.
A set $A$ equipped with a partial order $R$ is a \textbf{partially ordered set}
or \textbf{poset}, written $(A,R)$.
\end{definition}
\end{definitionbox}

\begin{remark*}[Standard quantified statement]
\[
\operatorname{PartialOrder}(R,A)
\Longleftrightarrow
\operatorname{ReflexiveRelation}(R,A)
\land
\operatorname{AntisymmetricRelation}(R,A)
\land
\operatorname{TransitiveRelation}(R,A).
\]
\end{remark*}

\begin{remark*}[Definition predicate reading]
$R$ partially orders $A$ exactly when it supports self-comparison,
antisymmetric comparison, and chaining.
\end{remark*}

\begin{remark*}[Negated quantified statement]
\[
\neg\operatorname{PartialOrder}(R,A)
\Longleftrightarrow
\neg\operatorname{ReflexiveRelation}(R,A)
\lor
\neg\operatorname{AntisymmetricRelation}(R,A)
\lor
\neg\operatorname{TransitiveRelation}(R,A).
\]
\end{remark*}

\begin{remark*}[Interpretation]
Partial orders allow incomparable elements while retaining enough structure for
comparison and chaining.
\end{remark*}

\begin{dependencies}
\begin{itemize}
  \item \hyperref[def:binary-relation]{Binary Relation}
  \item \hyperref[def:rel-reflexive]{Reflexive Relation}
  \item \hyperref[def:rel-antisymmetric]{Antisymmetric Relation}
  \item \hyperref[def:rel-transitive]{Transitive Relation}
\end{itemize}
\end{dependencies}

\begin{definition}[Strict Partial Order]
\label{def:algebraic-strict-partial-order}
A binary relation $<$ on $A$ is a \textbf{strict partial order} if it is
irreflexive, asymmetric, and transitive.
\end{definition}

\begin{remark*}[Standard quantified statement]
\[
\operatorname{StrictPartialOrder}(<,A)
\Longleftrightarrow
\operatorname{IrreflexiveRelation}(<,A)
\land
\operatorname{AsymmetricRelation}(<,A)
\land
\operatorname{TransitiveRelation}(<,A).
\]
\end{remark*}

\begin{remark*}[Definition predicate reading]
$<$ is a strict partial order exactly when it compares without self-comparison,
forbids reverse comparison, and chains transitively.
\end{remark*}

\begin{remark*}[Negated quantified statement]
\[
\neg\operatorname{StrictPartialOrder}(<,A)
\Longleftrightarrow
\neg\operatorname{IrreflexiveRelation}(<,A)
\lor
\neg\operatorname{AsymmetricRelation}(<,A)
\lor
\neg\operatorname{TransitiveRelation}(<,A).
\]
\end{remark*}

\begin{remark*}[Interpretation]
Strict partial orders encode comparison without equality cases.
\end{remark*}

\begin{dependencies}
\begin{itemize}
  \item \hyperref[def:rel-irreflexive]{Irreflexive Relation}
  \item \hyperref[def:rel-asymmetric]{Asymmetric Relation}
  \item \hyperref[def:rel-transitive]{Transitive Relation}
\end{itemize}
\end{dependencies}

\begin{definitionbox}{Definition (Total Order)}
\begin{definition}[Total Order]
\label{def:algebraic-total-order}
A partial order $\leq$ on $A$ is a \textbf{total order} if
\[
\forall a,b \in A,\quad a \leq b \lor b \leq a.
\]
A set equipped with a total order is a \textbf{totally ordered set} or
\textbf{chain}.
\end{definition}
\end{definitionbox}

\begin{remark*}[Standard quantified statement]
\[
\operatorname{TotalOrder}(\leq,A)
\Longleftrightarrow
\operatorname{PartialOrder}(\leq,A)
\land
\forall a,b\in A,\; a\leq b \lor b\leq a.
\]
\end{remark*}

\begin{remark*}[Definition predicate reading]
$\leq$ totally orders $A$ exactly when it is a partial order and every pair of
elements is comparable.
\end{remark*}

\begin{remark*}[Negated quantified statement]
\[
\neg\operatorname{TotalOrder}(\leq,A)
\Longleftrightarrow
\neg\operatorname{PartialOrder}(\leq,A)
\lor
\exists a,b\in A,\; \lnot(a\leq b)\land\lnot(b\leq a).
\]
\end{remark*}

\begin{remark*}[Interpretation]
Partial orders allow incomparable elements.  Total orders require every pair
to be comparable.  A strict order can be recovered from a non-strict order by
declaring $a<b$ exactly when $a\leq b$ and $a\neq b$.
\end{remark*}

\begin{dependencies}
\begin{itemize}
  \item \hyperref[def:algebraic-partial-order]{Partial Order}
  \item \hyperref[def:rel-total]{Total Relation}
\end{itemize}
\end{dependencies}
